\section{Conditional global probabilities: implementation and diagnostics}
\label{app:v504-global-study}
This appendix collects the detailed checks supporting the method in
Section~\ref{sec:v5-global} and its results in
Section~\ref{sec:v5-global-results}. The source and injection diagnostics
were computed before the 2015 extension; their original reference states
and sample sizes are retained. The new full-range probability curves are
Figures~\ref{fig:v5-global-2015} and \ref{fig:v5-global-combined}.
\subsection{Comparison with the paper and sigcorr}
\paragraph{Check against the paper and sigcorr.}
The official implementation~\cite{SigcorrPinned2023}, pinned at commit
\texttt{a66c08771a81}, reproduces our 2016 and combined correlation matrices
for both fit methods: the largest absolute difference is
$3.34\times10^{-16}$, and its Gaussian sampling factors reconstruct those
matrices within $4.89\times10^{-15}$. This verifies the covariance
normalization and sampling algebra. There is a substantive difference in
interpretation, however. The paper~\cite{AnanievRead2023} assumes a signed
likelihood-root field with null mean zero and unit variance. We supply
centered response directions $D$, retain a nonzero stress offset $a_m$ and
width $s_m$, and require a positive raw fit. These are explicit extensions,
not features that the reference package validates for HPS. Feeding the raw
biased stress scans into its zero-mean covariance formula instead changes
the profiled correlations by an RMS of 0.945 for 2016 and 0.953 for the
combination. Simply removing the centering would therefore be incorrect.
The calculations below are conditional stress-reference probabilities;
their normal quantiles are not calibrated particle-discovery significances.

\subsection{Large stress-reference offsets}
\paragraph{The large 2016 response and the 76 MeV comparison.}
The 2016 observed $r$ ranges from $-4.50$ to $3.42$, whereas its stress
response ranges from $-17.60$ to $13.71$. The large gold excursions therefore
describe the reference construction, not giant observed local signals.
The reference blends a low-mass fit and a broad continuation over
75--85~MeV; the latter retains a failed convergence flag and the construction
has development-sample overlap. The independent 1,000-scan sample reproduces
its bulk response: standardized means have magnitude below 0.099 and widths
range from 0.947 to 1.047, with no Holm-adjusted marginal-normality flags.
However, the nominal full-maximum comparison has $p=0.034$ and its rare tails
remain unvalidated. This is evidence that the approximation describes much
of this particular stress ensemble, not that the 2016 physical background
has been adequately qualified. A replacement reference would need independent
background controls and renewed validation, rather than selection because
it gives a more favorable observed significance.

At combined 76~MeV, $r=0.166$ gives nominal local $p_0=0.434$.
The stress response is $a=-8.700$ with width $s=0.979$, so
$(r-a)/s=9.053$. That number measures disagreement with a reference that
usually produces a large fitted deficit; it is not a $9\sigma$ resonance.
For 2016 alone at the same mass, $r=-2.460$ and $a=-14.653$. Although the
centered coordinate is 12.437, the negative fitted signal fails the declared
positive-fit gate and its conditional discovery probability is one.
The combined fit is slightly positive and passes that gate. This explains
both the apparent contrast between the plots and the sharp probability
jumps at sign changes. Under the separate raw-maximum ordering, every
archived combined simulated scan exceeds the observed maximum; the two orderings
cannot be exchanged after seeing their results.

\subsection{Resolving empty Gaussian tails}
\paragraph{Resolving the empty Gaussian tails.}
The original local plotting floor and ordinary Monte Carlo resolution were
separate issues. Removing the floor reveals all finite local tails.
The 2,000,000-field continuation still gives zero counts at six 2016 masses
(42--44, 55--56 and 71~MeV) and combined 76--77~MeV. Those ordinary-sampling
bounds remain in the ledger. This revision evaluates the same Gaussian
probabilities by conditional-mixture importance sampling
\cite{OwenMaximovChertkov2019}; it does not extrapolate a fitted tail or
change the reference to obtain smaller probabilities.

For an observed threshold $u$, the positive-fit gate makes the event
\begin{equation}
 \{T^*\geq u\}=\bigcup_i\{z_i^*\geq b_i\},\qquad
 b_i=\max(u,-a_i/s_i),\qquad Q=\sum_i[1-\Phi(b_i)].
 \label{eq:v502-tail-union}
\end{equation}
Choose a coordinate with probability $[1-\Phi(b_i)]/Q$ and draw a complete
correlated Gaussian field conditional on that coordinate exceeding its
threshold. If the draw exceeds $c$ coordinate thresholds, its weight is
$Q/c$. The mean weight is an unbiased estimator of the union probability.
The full covariance and scan grid are retained, including dataset-membership
boundaries. This additional integration method is separate from
\texttt{sigcorr}'s ordinary sampling and smooth-field upcrossing routines.
The latter cannot be applied unchanged to this varying gate threshold.

The fixed budget is two independent samples of 20,000 fields at each
positive-fit coordinate with standardized threshold at least 3.5:
23 coordinates in total, including every formerly empty Gaussian tail.
Checks cover independent and singular correlated fields with known
probabilities, plus eight overlap thresholds against the existing
2,000,000-field samples. The largest overlap discrepancy is 1.71 combined
standard errors; independent targeted replicates differ by at most 2.11.
The complete targeted run, including the overlap and basic analytic tests,
took 3.1 seconds with one worker, one linear-algebra thread and 207~MiB
peak resident memory. Relative Monte Carlo standard errors are below
0.33\%. Table~\ref{tab:v502-continuation} shows the newly resolved points.
The plotted intervals are pointwise approximate 95\% Monte Carlo intervals,
not simultaneous intervals or uncertainty in the Gaussian/background model.
The direct Poisson rare tails remain unresolved. A precisely integrated
biased Gaussian reference is still not a validated discovery probability.

\begin{table}[htbp]\centering\small
\begin{tabular}{lrrrr}\toprule
Scan & Mass [MeV] & Conditional global $p$ & $Z$ & Relative MC SE\\\midrule
2016 & 42 & $6.096\times10^{-20}$ & 9.067 & 0.16\%\\
2016 & 43 & $1.863\times10^{-18}$ & 8.687 & 0.17\%\\
2016 & 44 & $3.067\times10^{-9}$ & 5.813 & 0.23\%\\
2016 & 55 & $4.631\times10^{-11}$ & 6.479 & 0.21\%\\
2016 & 56 & $5.733\times10^{-13}$ & 7.112 & 0.20\%\\
2016 & 71 & $1.321\times10^{-7}$ & 5.147 & 0.24\%\\
Combined & 76 & $1.348\times10^{-17}$ & 8.459 & 0.15\%\\
Combined & 77 & $3.086\times10^{-9}$ & 5.812 & 0.20\%\\
\bottomrule\end{tabular}

\caption{Previously zero-count Gaussian tails resolved by conditional-mixture
sampling of the unchanged finite-grid stress model. Each estimate uses
40,000 weighted fields. The last column gives the relative Monte Carlo
standard error. The $Z$ column is a normal-quantile representation of the
conditional probability, not a particle significance. Ordinary Gaussian
sampling gives 0/2,000,000 at every listed point; the independent direct
Poisson study remains limited to 1,000 experiments.}
\label{tab:v502-continuation}
\end{table}

\fig{0.99}{../figures/v502_stress_reference_distributions.pdf}{Representative
background-only distributions from the saved 1,000 complete Poisson scans.
Green histograms show the signed fitted coordinate, gold curves the
unfluctuated-response Gaussian approximation, and black lines the observed
fit. These panels show why the apparently large centered scores at 2016
42~MeV and combined 76~MeV can coexist with nearly zero raw observed fits.
At 2016 76~MeV the observed fit is negative and is excluded by the
positive-fit gate. Agreement of the histograms with the gold curves tests
the approximation under that reference, not the physical correctness of
the reference spectrum.}
{fig:v502-stress-distributions}

\subsection{Fixed-mass injection recovery at 76 MeV}
\paragraph{Injected signals: what is already tested.}
Signal injection is built into the separate fixed-mass validation study in
Appendix~\ref{v5-sec:calibration}. Figure~\ref{fig:v502-injection-recovery}
replots its saved 76~MeV distributions for both 2016 and the combined fit,
using 500 independent Poisson toys per curve. The injected yield is fixed at
$A/\sigma_{\rm ref}=0,2,5$, where $\sigma_{\rm ref}$ is the frozen reference
uncertainty; these labels are injection sizes, not achieved significances.
For the combined local-GP control, the calibrated local $p<0.05$ criterion
is met in 23/500, 321/500 and 499/500 experiments, respectively. The
corresponding 2016 counts are 18/500, 311/500 and 500/500. Thus the fit
responds to an injected signal under that control, with residual offsets
visible in the distributions.

The stress-background panels are equally important. Their zero-injection
mean fitted yields are $-14.72\sigma_{\rm ref}$ for 2016 and
$-8.68\sigma_{\rm ref}$ for the combination. Adding signal moves the
fitted distribution to the right, but even the size-five injection leaves
it negative; all three strengths have 0/500 calibrated local rejections
under this stress construction. Calibration alone therefore does not repair
signal recovery under an unsuitable reference. These are local tests of
the two-truth calibration envelope, not a demonstration of full-scan
discovery power for the significance GP.

The practical gain of the significance GP is to approximate many correlated
background scans cheaply once the response has been evaluated
\cite{AnanievRead2023}. Compared with reading a nominal normal-tail p-value,
this exposes the effects of reference offsets, correlations and the search
range. It does not improve the fitted signal automatically or establish
superiority to a traditional background fit. A discovery-power comparison
would inject the same physical signal into coherent full spectra and apply
each complete search at the same validated global false-positive rate.
The present results support that next test while also identifying the
2016 reference problem that must be addressed first.

\fig{0.98}{../figures/v502_injection_recovery_76.pdf}{Saved 76~MeV signal-injection
experiments, with identical axes for the local-GP controls (top) and archived
stress backgrounds (bottom). Each curve has 500 independent Poisson spectra;
dotted vertical guides mark the true injected yields $0,2,5$ in units of
that scope's fixed $\sigma_{\rm ref}$. The distributions move with injection,
but the stress-background bias persists. These fitted-yield distributions
show recovery and its failure mode explicitly; injection size five does
not mean a measured $5\sigma$ signal. No new injection toys were generated
for this display.}
{fig:v502-injection-recovery}
\FloatBarrier

\subsection{Can the 2016 reference be improved?}
\paragraph{Can the 2016 reference be improved by repairing the blend?}
Figure~\ref{fig:v502-stress-source-controls} compares three predetermined
background constructions over all 142 frozen hypotheses: the archived
reference, its broad component alone with the same total normalization,
and a version integrating the blend weights inside each source bin.
No observed p-value enters either construction. The low-mass component is
a degree-five logistic--Chebyshev fit over 26--80~MeV; the 75--85~MeV
transition therefore uses some extrapolation. The broad component has a
failed source-fit convergence flag and threshold zeros. These are diagnostic
controls, not qualified candidate nulls.

Removing the low-mass component changes the 76~MeV response from
$-14.653$ to $-5.926$, and the 66~MeV response from $12.469$ to $-0.124$.
But at 42~MeV it worsens from $-9.215$ to $-23.524$; the full-grid RMS
increases from 4.967 to 9.430. More accurate integration changes the
response by at most 0.001135, ruling out midpoint blend weighting as the
cause of the large excursions in this paired comparison. The original
source is reproduced to floating-point precision, and exact scalar/batch
fits agree within $4.0\times10^{-12}$ in signed root. The final exact
replay takes 36 seconds and 364~MiB with one worker and no new Poisson
toys. Thus the source mixture matters, but neither tested modification
repairs 2016 across its search range. The existing response is retained;
qualifying a replacement needs background controls and signal-recovery
validation, not a cosmetically smoother curve.

\fig{0.98}{../figures/v502_stress_source_controls.pdf}{Controlled 2016
background-source comparison. Top: expected source spectra with identical
total normalization; gold is the archived blend and blue its broad component
alone. Middle: their exact unfluctuated signed responses, using the same
frozen kernels. Bottom: the difference in deterministic signed response,
$r_{\rm integrated}-r_{\rm midpoint}$, when the blend weights are integrated
inside each source bin. This is neither an observed mass residual nor a
fitted signal: both curves come from smooth, signal-free expected spectra.
The maximum absolute change is only 0.001135 in $r$. A localized smooth
shape on this very small vertical scale therefore does not show extraction
of a perfect signal. Gray shading marks 75--85~MeV. The broad-only
control improves some masses but makes the full-grid response worse; it
retains a failed source-fit flag and is not adopted as a replacement null.
No new Poisson or injected-signal toys are used here.}
{fig:v502-stress-source-controls}
\FloatBarrier

\subsection{Covariance and scan-grid checks}

\paragraph{Covariance and injection follow-up (v5.1.1).}
The new study tests two links in the approximation: whether the response
matrix describes the dependence between fitted masses, and whether a known
signal produces the expected local recovery and neighboring response.
The observed spectra and original study products were unchanged by that support study.
These are conditional diagnostics of the fixed analysis and generating
spectra, not a new global calibration.

For the extended-grid covariance check, the same first 256 archived
background-only Poisson spectra are evaluated at every coordinate and
centered by their empirical means. Their covariance
and correlation are compared with $D^\top D$ and $R$ from
Eqs.~\ref{v5-eq:field-response}--\ref{eq:v501-field-correlation}.
Bootstrap resampling keeps each whole scan together, preserving the
correlations among masses. The resulting uncertainties quantify Monte Carlo
precision; they do not include uncertainty in the stress-background model.
The eigenmodes describe correlated patterns in the scan. An effective
mode count is a covariance summary, not a calibrated number of search trials.
The RMS difference between empirical and response-predicted correlations
is 0.0556 for 2015, 0.0626 for 2016 and 0.0594 for 2021. The denser
2016 grid resolves the shape of the stress response, whose signed root
still spans $-17.86$ to $+13.71$; grid refinement does not repair that
background offset.

The 2015 pilot covers 19--100~MeV at integer masses. Above 90~MeV the
kernel coordinates are held at their 90~MeV values; each added mass still
has its own recomputed fit and source-bin response. The 2016 grid retains
39--180~MeV and adds half-MeV points from 70.5 to 80.5~MeV, with log-linear
interpolation of the saved kernel coordinates. The released 2021 sample
retains 50--250~MeV at integer masses. The grids therefore contain
82, 153 and 201 points. Re-evaluating the same 256 complete spectra at
all coordinates provides the new covariance blocks; unrelated local toys
are not joined. The earlier 1,000-scan native-grid benchmark is retained
separately. A timing-based amendment fixed this reduced covariance sample
and grid before inspecting their results. That study alone did not enlarge the released observed search. Version 5.0.4
now uses its endpoint prescription and response matrix, together with the
new conditional limits in Appendix~\ref{app:v504-extension}.

\fig{0.99}{../figures/v511_covariance_comparison.pdf}{Response-predicted
correlations and their empirical comparison using the same 256 complete
background-only Poisson spectra at all coordinates. Each row uses its
actual mass grid; difference panels have separate color scales. The
original 1,000-scan native-grid result is retained separately and is not
spliced into these extended matrices. This checks central dependence
under the specified stress reference, not the rare global tail.}
{fig:v511-covariance}
\FloatBarrier

\subsection{Paired Poisson injection distributions}
\fig{0.99}{../figures/v511_injection_distributions.pdf}{Fixed-mass Poisson
injection examples under local GP controls (top) and the archived stress
truth (bottom). Each curve uses 64 experiments, paired across the three
strengths within a mass and truth. Dotted lines show fits to unfluctuated
expectations. The horizontal axis is the signed fitted root; the labels
$A/\sigma_{\rm ref}=0,2,5$ specify injected yields, not achieved
significances. The same reference uncertainty is used for both truths
at a fixed mass. Persistent stress offsets are especially clear at
2016 76~MeV. No global discovery efficiency is assigned.}
{fig:v511-injections}
\FloatBarrier

The fixed-mass Poisson experiments measure conditional signal recovery
at their declared anchors. They do not scan each fluctuated injection
across the full mass range, and no new global false-positive threshold or
discovery power is inferred. The significance GP remains useful for
propagating a specified background response across the mass scan;
background qualification and global-tail validation remain separate
requirements.
\FloatBarrier

\subsection{Why the 76 MeV probability spike is not a narrow mass peak}
\label{app:v504-width-review}
The statistical review first distinguishes the dataset and the statistic.
At 76~MeV the 2016 fit has $r=-2.4599$, stress expectation
$a=-14.6532$, and response width $s=0.9804$. Its positive-fit gate is
closed. The combined fit has $r=0.1656$, $a=-8.7001$ and $s=0.9794$,
giving $(r-a)/s=9.0526$ with the gate open. At 75~MeV the combined
raw fit is $r=-0.6659$ and the gate is closed. This change in the gate,
followed by a rapid change in the reference offset, creates the sharp
feature in the plotted probability. The 2015 extension changes the global
maximum distribution but leaves these fixed-mass fits below 90~MeV intact.

The adopted 2016 detector width is $\sigma_m=3.0979$~MeV at 76~MeV,
or FWHM $=7.2949$~MeV for a Gaussian. Interpolating the old plotted global
$Z$ nodes gives an apparent half-height width of about 1.90~MeV. That is
a property of a transformed, gated scan, not a fitted invariant-mass width.
Even a known resolution-shaped injection has a different scan response:
the five-reference-error injection gives a positive $\Delta r$ lobe about
5.74~MeV wide. The training mask moves and the GP is retrained at each
hypothesis, so the scan acts as a filter and can develop negative neighbors.
Equation~\eqref{eq:v511-injection-projection} is the appropriate first-order
description of that response; the detector line shape is a different object.

All twelve saved fixed-width and GP-treatment variants examined at 76~MeV
retain negative 2016 fits and zero bounded signal. The saved resolution-nuisance
fit also gives zero positive signal. Its width coordinate at zero amplitude
cannot be read as a measured narrow width. The 2016 90--92~MeV region does
have positive raw fits, about $r=3.1$--$3.4$, which makes a mass-residual
inspection meaningful. This is a reason to investigate its shape and rate
agreement, not to assign it a physical origin. A width inference would require
a fixed-bin signal-plus-background model and an independently supported
resolution constraint, together with the effect of selecting the region.
The review adds no such width measurement or calibrated discovery claim.

\fig{1.0}{../figures/v510_histogram_binning.pdf}{Archived v5.1.0 histogram-binning
comparison using exact pointwise fits at the saved kernel settings. Histogram
width changes both the GP prediction and the signal fit. These individual
campaign scans are not a rebinned calibration of the combined positive-fit gate.}
{fig:v505-binning}

\subsection{Joint binning check around the 76 MeV gate crossing}
\label{app:v505-joint-binning}
The archived v5.1.0 scans above describe the individual campaigns. For this
revision, a short joint scan tests the effect on the combined signed fit
itself. Five binning choices are fixed before the refits: the nominal widths,
twice and four times those widths, and one additional origin for each coarse
choice shifted by half its bin width. The tested masses are the integers
from 72 through 80 MeV. The choice of this region follows the question about
the displayed spike; no new discovery test is assigned to the comparison.

The original analysis widths are 0.25 MeV for 2015 and 2016 and 0.625 MeV
for 2021. Counts are summed in adjacent original bins, with exact conservation
inside each included group. An incomplete group at the far upper training
edge is omitted, as in the original rebin operation; shifting the origin
also removes leading support bins. The effective endpoints are saved for
every campaign and choice. Thus this check includes the small change in
sideband sampling that accompanies the selected bin boundaries; it does not
isolate bin origin from support-edge effects.

Each tested mass retains its archived kernel amplitude and length scale.
The log-count GP is conditioned again on the rebinned sidebands, excluding
bin centers within $\pm2.25\sigma_m$. The count covariance and
bin-integrated Gaussian template are reconstructed, and the joint Poisson
likelihood is fitted with a common signed coupling and independent GP
background constraints. The native-bin prompt density, detector resolution
and radiative conversion are held fixed. The nominal 76 MeV fit reproduces
the reported $r=0.166$ to the quoted precision, and its individual roots
agree with the archived values. Every fitted expectation is positive and
all 180 fitted mass/scope/binning coordinates converge.

\begin{table}[htbp]\centering\small
\begin{tabular}{lrrr}\toprule
Widths relative to nominal & Origin shift & $r_{\rm comb}(76)$ & First $r>0$ [MeV]\\\midrule
$1\times$ & 0 & $+0.166$ & 76\\
$2\times$ & 0 & $-0.216$ & 77\\
$2\times$ & half a coarse bin & $+0.530$ & 76\\
$4\times$ & 0 & $-0.137$ & 77\\
$4\times$ & half a coarse bin & $-0.695$ & 77\\
\bottomrule\end{tabular}
\caption{Signed joint-profile response in the bounded v5.0.5 diagnostic.
``First'' refers only to the tested 72--80 MeV grid. The shifted origins
are applied separately using each campaign's bin width.}
\label{tab:v505-joint-binning}\end{table}

\fig{1.0}{../figures/v505_76_joint_binning.pdf}{Observed signed fit scans
under the five declared bin choices. The horizontal zero line is the
positive-fit gate boundary. At 76 MeV, the combined fit can lie on either
side of it while the 2016 component remains negative. Curves connect tested
integer masses; no sub-MeV zero-crossing precision is implied.}
{fig:v505-joint-binning}
The saved CSV and protocol record the bin widths, origin shifts, effective
supports and convergence diagnostics. This small comparison draws no toys,
reoptimizes no kernels and computes no global probabilities. A full binning
choice for unblinding still requires the corresponding closure and
calibration checks.
